\documentclass[11pt]{article}

\usepackage[T1]{fontenc}
\usepackage[utf8]{inputenc}
\usepackage{lmodern}
\usepackage{amsmath,amssymb,amsthm}
\usepackage{booktabs}
\usepackage[margin=1in]{geometry}
\usepackage{microtype}
\usepackage{url}
\usepackage{xcolor}
\usepackage[colorlinks=true,linkcolor=blue,citecolor=blue,
  urlcolor=blue,
  pdftitle={Coloop Multiplication for Radial Mobius Hilbert Series},
  pdfauthor={Xiaojian Zeng}]{hyperref}
\usepackage[backend=biber,style=alphabetic]{biblatex}
\addbibresource{references.bib}

\newtheorem{theorem}{Theorem}
\newtheorem{lemma}[theorem]{Lemma}
\newtheorem{proposition}[theorem]{Proposition}
\theoremstyle{remark}
\newtheorem{remark}[theorem]{Remark}

\newcommand{\Q}{\mathbb{Q}}
\newcommand{\cl}{\operatorname{cl}}

\title{Coloop Multiplication for Radial M\"obius Hilbert Series:\\
the Connected Graphic Classification Through Five Vertices}
\author{Xiaojian Zeng\\
PBC School of Finance, Tsinghua University, 100190\\
\texttt{xiaojianzeng@foxmail.com}}
\date{}

\begin{document}
\maketitle

\begin{abstract}
For the radial cyclic module of the graded M\"obius algebra, adjoining a
coloop suggests multiplication of the Hilbert series by $1+q$.  We give a
complete computer-assisted classification of this identity for cycle
matroids of connected simple graphs on at most five vertices.  The identity
holds through order four.  At order five it fails exactly when the graph
complement is a linear forest.  We determine all seven nonzero differences;
the unique failure with the fewest edges is
$K_5-P_5\cong\Theta(1,2,3)$, whose Hilbert vectors are
$(1,6,11,6,1)$ and $(1,7,19,19,7,1)$ before and after adjoining the
coloop.  The proof exhausts all $31$ graph isomorphism classes, constructs
every flat and cyclic space over $\Q$, and is checked by a separate exact
implementation based on connected set partitions.  The result is confined
to this finite graphic range and does not settle the general operations
question for the radial cyclic module.
\end{abstract}

\section{Context and statement}

Cycle matroids are among the motivating examples in Whitney's introduction
of matroids \cite{Whitney1935}.  The graded M\"obius algebra of a matroid is
also a natural module-theoretic object in modern combinatorial Hodge theory
\cite{BradenEtAl2026}.  Sinclair \cite{Sinclair2026} introduced a canonical
radial lowering operator on this algebra and the Hilbert series of the cyclic
module generated jointly by it and the atom lowerings.  He asked how that
module behaves under direct sum, deletion, contraction, and duality, while
noting that simple formulas for the Hilbert series fail
\cite[Question~6.3]{Sinclair2026}.

This note isolates the direct sum with one coloop in the smallest connected
graphic range.  Relative to \cite{Sinclair2026}, the contribution is the
exact classification through five vertices, including the first failing
order in that range, the least-edge uniqueness statement, and every nonzero
difference.  Sinclair's Example~3.2 supplies the value for $K_5$ but not its
coloop extension or the classification below.  No assertion is made here
about graphs of larger order or about the full operations question.

We now fix the conventions.  Let $G=(V,E)$ be a connected simple graph.  For
$A\subseteq E$, let $c(V,A)$ count all components of the spanning subgraph
$(V,A)$, isolated vertices included, and set
\[
  r(A)=|V|-c(V,A).
\]
The cycle-matroid closure is
\[
  \cl_G(A)=A\cup
  \{uv\in E:u\text{ and }v\text{ are connected in }(V,A)\}.
\]
Thus $X\subseteq E$ is a flat when $\cl_G(X)=X$.  Write
$X\prec Y$ when $X$ is covered by $Y$ in the flat lattice.

Work over $\Q$ with basis $\{g_X:X\text{ is a flat}\}$.  For $e\in E$,
define the lowering operators
\begin{align*}
 L_e g_Y
   &=\sum_{\substack{X\prec Y\\e\in Y\setminus X}}g_X,\\
 Dg_Y
   &=r(Y)\sum_{X\prec Y}
       \frac{|Y\setminus X|}{|E\setminus X|}g_X.
\end{align*}
These are the rescaled flat-cover formulas of
\cite[Proposition~2.4]{Sinclair2026}, specialized to a simple graphic
matroid.  A denominator is never zero, because a lower member of a cover is
not the top flat.  Starting at the top flat, put
\[
 C_0(M)=\Q g_E,
 \qquad
 C_{j+1}(M)=D C_j(M)+\sum_{e\in E}L_e C_j(M),
\]
and
\[
 H_M(q)=\sum_{j=0}^{r(M)}\dim_{\Q}C_j(M)q^j.
\]
If $c$ is a new coloop, let $M^+=M\oplus U_{1,1}$ and define
\[
  \Delta_G(q)=H_{M^+}(q)-(1+q)H_M(q).
\]
Graphically, $M^+$ is realized by attaching a new leaf edge to $G$.  For the
one-vertex graph $E=\varnothing$, the conventions give $H_M=1$.

A \emph{linear forest} is a disjoint union of paths; an isolated vertex is a
copy of $P_1$.

\begin{theorem}[Complete classification through five vertices]
\label{thm:classification}
Let $G$ be a connected simple graph with $1\leq |V(G)|\leq5$ and let all
objects be as above.
\begin{enumerate}
\item If $|V(G)|\leq4$, then $\Delta_G(q)=0$.
\item If $|V(G)|=5$, then $\Delta_G(q)\neq0$ if and only if the complement
  $\overline G$ is a linear forest.
\item The seven nonzero cases, including the two Hilbert vectors that certify
  each difference, are listed in Table~\ref{tab:nonzero}.
\item Within this range, order five is the first failing order.  Among the
  five-vertex failures, the minimum number of edges is six, and the unique
  such graph up to isomorphism is
  $K_5-P_5\cong\Theta(1,2,3)$.  Its two Hilbert vectors are
  $(1,6,11,6,1)$ and $(1,7,19,19,7,1)$.
\end{enumerate}
\end{theorem}

\begin{table}[htbp]
\centering
\small
\begin{tabular}{@{}lcll@{}}
\toprule
$\overline G$ & $|E(G)|$ & $(\dim C_j(M))_j$ &
  $(\dim C_j(M^+))_j$ \\
\midrule
$P_5$                  & 6  & $(1,6,11,6,1)$ & $(1,7,19,19,7,1)$ \\
$P_4\sqcup P_1$        & 7  & $(1,7,13,7,1)$ & $(1,8,22,22,8,1)$ \\
$P_3\sqcup P_2$        & 7  & $(1,8,15,7,1)$ & $(1,9,25,24,8,1)$ \\
$P_3\sqcup2P_1$        & 8  & $(1,9,16,8,1)$ & $(1,10,26,26,9,1)$ \\
$2P_2\sqcup P_1$       & 8  & $(1,9,20,8,1)$ & $(1,10,33,28,9,1)$ \\
$P_2\sqcup3P_1$        & 9  & $(1,10,22,9,1)$ & $(1,11,36,31,10,1)$ \\
$5P_1$                  & 10 & $(1,11,25,10,1)$ & $(1,12,40,35,11,1)$ \\
\bottomrule
\end{tabular}

\medskip
\begin{tabular}{@{}lc@{}}
\toprule
$\overline G$ & $\Delta_G(q)$ \\
\midrule
$P_5,\ P_4\sqcup P_1,\ P_3\sqcup P_2$ & $2q^2+2q^3$ \\
$P_3\sqcup2P_1$                         & $q^2+2q^3$ \\
$2P_2\sqcup P_1,\ P_2\sqcup3P_1,\ 5P_1$ & $4q^2$ \\
\bottomrule
\end{tabular}
\caption{All nonzero differences at order five.}
\label{tab:nonzero}
\end{table}
\section{Two exact models for graphic flats}

The proof is finite, but its mathematical specification does not depend on a
particular graph or flat encoding.  The generating computation uses closure
of edge subsets.  A separate checker uses the following partition model.

\begin{lemma}[Connected-partition model]
\label{lem:partition}
The flats of $M(G)$ are in bijection with the set partitions $\pi$ of $V$
whose blocks induce connected subgraphs of $G$.  The flat associated with
$\pi$ is
\[
 X(\pi)=\{uv\in E:u\text{ and }v\text{ lie in the same block of }\pi\},
\]
and its rank is $|V|-|\pi|$.
\end{lemma}

\begin{proof}
Let $X$ be a flat, and let $\pi$ be the component partition of $(V,X)$.
Every block is connected in $G$.  Closure adds precisely all edges of $G$
whose endpoints are in one component of $(V,X)$.  Since $X$ is closed, it
already contains all of them, so $X=X(\pi)$.

Conversely, suppose every block of $\pi$ induces a connected subgraph of
$G$.  Then the components of $(V,X(\pi))$ are exactly the blocks of $\pi$.
Every graph edge internal to one component already belongs to $X(\pi)$, so
$X(\pi)$ is closed.  The rank formula follows from the number of spanning
components.
\end{proof}

If $X\subsetneq Y$ are flats of adjacent ranks, they form a cover; therefore
either implementation can determine the exact flat-cover relation.  On a
ranked flat layer, the two matrices used in the computation have entries
\begin{align}
 (L_e)_{X,Y}&=\mathbf{1}_{\{e\in Y\setminus X\}},
 \label{eq:Le}\tag{1}\\
 D_{X,Y}&=r(Y)\frac{|Y\setminus X|}{|E|-|X|}
 \label{eq:D}\tag{2}
\end{align}
when $X\prec Y$, and zero otherwise.  Thus the definitions are reduced to
finite set operations and exact rational arithmetic.

\section{Exhaustion and rational-rank certificates}

We describe enough of the certificate construction to make clear why its
successful replay proves Theorem~\ref{thm:classification}.

\subsection{All graph isomorphism classes}

For each $n=1,\ldots,5$, the generator visits all
$2^{\binom n2}$ labeled edge masks, rejects the disconnected masks, and
minimizes the adjacency word over all $n!$ vertex permutations.  It also
computes the automorphism group size of each representative.  The resulting
counts are shown in Table~\ref{tab:coverage}.  The orbit-size sum agrees with
the independently counted connected labeled graphs in every order, so the
$31$ retained representatives cover the quantified domain exactly once up
to isomorphism.

\begin{table}[htbp]
\centering
\begin{tabular}{@{}rrrrr@{}}
\toprule
$n$ & labeled masks & connected labeled & representatives & orbit-size sum \\
\midrule
1 & 1    & 1   & 1  & 1   \\
2 & 2    & 1   & 1  & 1   \\
3 & 8    & 4   & 2  & 4   \\
4 & 64   & 38  & 6  & 38  \\
5 & 1024 & 728 & 21 & 728 \\
\bottomrule
\end{tabular}
\caption{Exact graph-orbit coverage.}
\label{tab:coverage}
\end{table}

Relabeling vertices permutes the edges and flats.  The induced map
$g_X\mapsto g_{\sigma X}$ intertwines $D$ and permutes the family
$(L_e)_{e\in E}$.  It therefore carries every $C_j(M(G))$ isomorphically to
the corresponding space for $\sigma G$.  One representative per orbit is
enough.

\subsection{Flat and cyclic-space certificates}

For a representative with $m$ edges, the generator visits all $2^m$ edge
subsets.  Union--find computes the components of $(V,A)$ and hence
\[
 \cl_G(A)=\{uv\in E:u,v\text{ have the same component root in }(V,A)\}.
\]
Its fixed points give every flat.  Covers are then found by literal
intermediate-flat exclusion.  The program asserts, for every edge subset,
extensivity, idempotence, and rank preservation of closure, and also checks
that literal covers agree with inclusions between adjacent rank layers.

Starting with the unit vector at the top flat, assume that an exact row basis
for $C_j$ has been obtained in one flat layer.  Apply every matrix
\eqref{eq:Le} and the matrix \eqref{eq:D} to every basis vector, concatenate
the images, and perform Gaussian elimination over $\Q$.  The resulting rows
span
\[
  DC_j+\sum_{e\in E}L_eC_j=C_{j+1}.
\]
Induction on $j$ proves that their exact ranks are the coefficients of
$H_M$.  The full certificate serializes every flat layer, coordinate order,
and sparse rational echelon basis, rather than only the final dimensions.

For $M^+$ the computation attaches a new vertex by one pendant edge and
repeats the entire construction.  The pendant edge is a bridge, hence a
coloop, so the resulting cycle matroid is $M(G)\oplus U_{1,1}$.  No direct-sum
formula is assumed in calculating $H_{M^+}$.

The separate replay checker does not call the subset-closure implementation.
It generates the flats from Lemma~\ref{lem:partition}, constructs covers and
the matrices \eqref{eq:Le}--\eqref{eq:D} afresh, and computes exact column
spaces with SymPy \cite{MeurerEtAl2017}.  For all $31$ graphs and all $31$
pendant-edge extensions, its flat layers equal the serialized layers; each
serialized basis is linearly independent and spans the freshly recomputed
cyclic space in every degree.  A separate execution of the generator also
reproduces all $62$ certificates exactly.  This certificate pattern is
related in spirit to exhaustive graph-isomorphism and exact-arithmetic checks
of algebraic multiplication maps in \cite{Zeng2026Whiskered}; here all
claimed ranks are over $\Q$, and no floating-point or modular rank is used.

The replay returns zero differences for all ten representatives of orders at
most four and precisely the seven rows of
Table~\ref{tab:nonzero} among the $21$ representatives of order five.  The
$K_5$ row is $(1,11,25,10,1)$, agreeing with
\cite[Example~3.2]{Sinclair2026}.  These exact outputs prove the first three
parts of Theorem~\ref{thm:classification}: a linear forest on five vertices
is determined up to isomorphism by the path sizes in an integer partition of
five, and the seven partitions are exactly the seven complement types in the
table.

\subsection{Least-edge uniqueness}

It remains only to extract the final structural conclusion from the table.
A linear forest on five vertices with $k$ components has $5-k$ edges.  Among
the seven possible types, only $P_5$ has four edges.  Hence only its
complement has $10-4=6$ edges; all other failures have at least seven edges.
Thus the unique least-edge failure is $K_5-P_5$.

For completeness, take vertices $0,1,2,3,4$.  The theta graph formed by the
three internally disjoint paths
\[
  0-1,\qquad 0-2-1,\qquad 0-3-4-1
\]
has complement path $0-4-2-3-1$.  It is therefore
$K_5-P_5\cong\Theta(1,2,3)$.  Its two Hilbert vectors are the first row of
Table~\ref{tab:nonzero}, completing the proof of
Theorem~\ref{thm:classification}.

\section{Reproducibility, assistance, and limits}

The ancillary archive contains the exact generator
\path{independent_subset_check.py}, its complete JSON certificate
\path{subset_closure_census.json}, the connected-partition replay
\path{referee_verify.py}, and the detailed replay output
\path{referee_replay.json}.  The certificate can be regenerated with
Python~3 using only standard-library rational arithmetic; the separate replay
uses SymPy~1.14.0.  These artifacts cover exactly the finite domain stated in
Theorem~\ref{thm:classification}.

OpenAI Codex was used in a bounded research workflow to help generate and
inspect code, execute the exact computations, conduct a separate replay of
the certificate, check bibliographic records, and draft this manuscript.
Language-model output is not used as mathematical evidence: the finite
claims are supported by inspectable source code, complete exact certificates,
and a separate exact implementation.  The author is responsible for the
statement and the manuscript.

Nothing here addresses connected graphs on six or more vertices,
disconnected or nonsimple graphs, arbitrary matroids, or a general formula
for direct sum, deletion, contraction, or duality.  In particular, the result
does not answer Sinclair's Question~6.3.  The literature comparison in this
note is deliberately limited to the cited records and is not a claim of
priority.

\printbibliography

\end{document}
